% ============================================================
% Chapter: Numerical Modelling and Dynamic Assessment
% Assessed MLOs: MLO-2
% Excellent-grade rubric criteria:
% - Numerical modelling approach (5%): the conceptual model is appropriate,
%   simplified deliberately, and justified against the physics of the problem.
% - Numerical analysis (5%): the equations of motion, solution strategy, and
%   implementation choices are coherent and suited to offshore loading.
% - Numerical and advanced dynamic assessment (5%): results are validated or
%   benchmarked, interpreted critically, and used to delimit applicability.
% ============================================================

\section{Numerical Modelling and Dynamic Assessment}
\pagecheck{8--9}

% Rubric Checklist: Numerical modelling approach (5%); Numerical analysis (5%); Numerical and advanced dynamic assessment (5%).
\instructionplaceholder{Unified chapter logic}{This chapter combines the old modelling approach, numerical analysis, and advanced dynamic assessment into one coherent workflow. The reader should be able to follow the path from conceptual model, to numerical implementation, to engineering interpretation without jumping across chapters.}

\subsection{Conceptual model definition}
\instructionplaceholder{Required evidence}{Describe the structural idealisation used for the problem. State the bodies, beams, rods, springs, masses, damping terms, hydrodynamic representations, and boundary conditions included in the model. Explain why this abstraction captures the physics that matter for the design question.}

\subsection{Governing equations and numerical method}
\instructionplaceholder{Required evidence}{Present the equations of motion or the governing numerical formulation at a level appropriate for a technical report. State which solution method is used, how loading is applied, and what assumptions are made about linearity, coupling, damping, and time or frequency domain treatment.}

\subsection{Implementation and verification}
\instructionplaceholder{Required evidence}{Explain how the method is implemented in Python, Matlab, or another platform. State what was checked to verify the implementation, such as unit checks, limiting-case comparisons, convergence studies, or comparison to analytical solutions. The code itself can be moved to an appendix.}

\subsection{Dynamic response assessment}
\instructionplaceholder{Required evidence}{Report the dynamic response results needed for design judgement. These may include motions, accelerations, internal forces, mooring loads, modal properties, hydroelastic response, or other quantities relevant to the concept. Interpret the results, identify governing cases, and state clearly what the model can and cannot be trusted to predict.}
